\subsection{The Section-Average Trace}\label{sec:trace-cowper0}

The section-average trace is defined by the $L^2$-optimal rigid approximation 
of \cref{prop:align-cowper}:
\begin{equation}
\bar{\bm{u}}(x) := \SectionAverage{\hat{\bm{u}}} - \bar{\bm{\theta}} \times \CentroidLocation, 
\qquad 
\bar{\bm{\theta}}(x) := \RotationAverage{\hat{\bm{u}}},
% \label{eq:def-section-avg-trace}
\end{equation}
where \(\SectionAverage{\cdot}\) and \(\RotationAverage{\cdot}\) are the functionals defined by \eqref{eq:def-section-averages}. 
The warping amplitude $\bm{\alpha}(x)$ is determined by the residual:
\begin{equation}
\Phi(\bm{r})\,\bm{\alpha}(x) = \hat{\bm{u}}(x,\bm{r}) - \bar{\bm{u}}(x) - \bar{\bm{\theta}}(x) \times \bm{r}.
\end{equation}

\begin{remark}
The residual $\tilde{\bm{u}} := \hat{\bm{u}} - (\bar{\bm{u}} + \bar{\bm{\theta}} \times \bm{r})$ 
is $L^2$-orthogonal to the section-rigid modes:
\begin{equation*}
\int_{\Section} \tilde{\bm{u}} \, \mathrm{d}A = \mathbf{0}, 
\qquad 
\int_{\Section} \bm{r} \times \tilde{\bm{u}} \, \mathrm{d}A = \mathbf{0}.
\end{equation*}
This is the defining property of the section-average trace: it extracts the 
section-rigid part that minimizes the $L^2$-norm of the warping residual.
\end{remark}

We now compute the trace matrices $\TraceDeDpRoot, \TraceDeDpOne$ by 
applying \eqref{eq:def-section-avg-trace} to each Ieşan mode systematically. 

\paragraph{Extension.}

The extension mode displacement is
\begin{equation*}
\hat{\bm{u}}^{(\varepsilon)} = (x\,\FrameBasis - \nu\,\bm{r})\,a_\varepsilon.
\end{equation*}
Applying the section-average operators:
\begin{align*}
\SectionAverage{\hat{\bm{u}}^{(\varepsilon)}} &= x\,\FrameBasis\,a_\varepsilon - \nu\,\CentroidLocation\,a_\varepsilon, \\
\RotationAverage{\hat{\bm{u}}^{(\varepsilon)}} &= \RotationAverage{x\,\FrameBasis - \nu\,\bm{r}}\,a_\varepsilon = \mathbf{0},
\end{align*}
since $\RotationAverage{\FrameBasis} = \mathbf{0}$ (constant axial field has no rotation content) 
and $\RotationAverage{\bm{r}} = \mathbf{0}$ (the field $\bm{r}$ is a pure translation by definition 
of centroidal coordinates). Thus:
\begin{equation*}
\bar{\bm{u}}^{(\varepsilon)}(x) = x\,\FrameBasis\,a_\varepsilon - \nu\,\CentroidLocation\,a_\varepsilon, 
\qquad 
\bar{\bm{\theta}}^{(\varepsilon)}(x) = \mathbf{0}.
\end{equation*}
The strains are:
\begin{equation*}
\bar{\bm{\gamma}}^{(\varepsilon)} = (\bar{\bm{u}}^{(\varepsilon)})' + \FrameBasis \times \bar{\bm{\theta}}^{(\varepsilon)} 
= \FrameBasis\,a_\varepsilon, 
\qquad 
\bar{\bm{\kappa}}^{(\varepsilon)} = (\bar{\bm{\theta}}^{(\varepsilon)})' = \mathbf{0}.
\end{equation*}
Hence $\bm{e}^{(\varepsilon)}(x) = (\FrameBasis, \mathbf{0})^\top a_\varepsilon$, which is constant in $x$.

\paragraph{Bending.}

The bending mode displacement is
\begin{equation*}
\hat{\bm{u}}^{(a)} = \Big( -\tfrac{1}{2}x^2\,\mathbf{1}_\Omega 
- \nu\,\mathrm{dev}(\bm{r} \otimes \bm{r}) 
+ x\,\FrameBasis \otimes \bm{r} \Big) \bm{a}_\Omega.
\end{equation*}
Applying the section-average operators:
\begin{align*}
\SectionAverage{\hat{\bm{u}}^{(a)}} &= -\tfrac{1}{2}x^2\,\bm{a}_\Omega 
- \nu\,\SectionAverage{\mathrm{dev}(\bm{r} \otimes \bm{r})}\,\bm{a}_\Omega 
+ x\,(\CentroidLocation \cdot \bm{a}_\Omega)\,\FrameBasis, \\
\RotationAverage{\hat{\bm{u}}^{(a)}} &= -x\,(\FrameBasis \times \bm{a}_\Omega) 
+ \RotationAverage{-\nu\,\mathrm{dev}(\bm{r} \otimes \bm{r})}\,\bm{a}_\Omega,
\end{align*}
where we used $\RotationAverage{(\FrameBasis\times\bm{a}_{\Section})\times \bm{r}} = -\,(\FrameBasis \times \bm{a}_\Omega)$. 
Then with \(\RotationAverage{\PlaneShape} \in \mathbb{R}^{3 \times 2}\):
\begin{align*}
\bar{\bm{\theta}}^{(a)}(x) &= -x\,(\FrameBasis \times \bm{a}_\Omega) + \RotationAverage{\PlaneShape}\,\bm{a}_\Omega, \\
\bar{\bm{u}}^{(a)}(x) &= -\tfrac{1}{2}x^2\,\bm{a}_\Omega 
- \nu\,\SectionAverage{\mathrm{dev}(\bm{r} \otimes \bm{r})}\,\bm{a}_\Omega 
+ x\,(\CentroidLocation \cdot \bm{a}_\Omega)\,\FrameBasis 
- \bar{\bm{\theta}}^{(a)} \times \CentroidLocation.
\end{align*}
The strains are:
\begin{align*}
\bar{\bm{\kappa}}^{(a)} &= (\bar{\bm{\theta}}^{(a)})' = -\FrameBasis \times \bm{a}_\Omega, \\
(\bar{\bm{u}}^{(a)})' &= -x\,\bm{a}_\Omega + (\CentroidLocation \cdot \bm{a}_\Omega)\,\FrameBasis 
- (\bar{\bm{\theta}}^{(a)})' \times \CentroidLocation \\
&= -x\,\bm{a}_\Omega + (\CentroidLocation \cdot \bm{a}_\Omega)\,\FrameBasis 
+ (\FrameBasis \times \bm{a}_\Omega) \times \CentroidLocation.
\end{align*}
Using the identity $(\FrameBasis \times \bm{a}_\Omega) \times \CentroidLocation 
= (\FrameBasis \cdot \CentroidLocation)\,\bm{a}_\Omega - (\bm{a}_\Omega \cdot \CentroidLocation)\,\FrameBasis 
= -(\CentroidLocation \cdot \bm{a}_\Omega)\,\FrameBasis$ (since $\CentroidLocation \perp \FrameBasis$):
\begin{equation*}
(\bar{\bm{u}}^{(a)})' = -x\,\bm{a}_\Omega.
\end{equation*}
Thus:
\begin{align*}
\bar{\bm{\gamma}}^{(a)} &= (\bar{\bm{u}}^{(a)})' + \FrameBasis \times \bar{\bm{\theta}}^{(a)} \\
&= -x\,\bm{a}_\Omega + \FrameBasis \times \big( -x\,(\FrameBasis \times \bm{a}_\Omega) + \RotationAverage{\PlaneShape}\,\bm{a}_\Omega \big) \\
&= -x\,\bm{a}_\Omega + x\,\bm{a}_\Omega + \FrameBasis \times \RotationAverage{\PlaneShape}\,\bm{a}_\Omega \\
&= (\FrameBasis \times \RotationAverage{\PlaneShape})\,\bm{a}_\Omega,
\end{align*}
where we used $\FrameBasis \times (\FrameBasis \times \bm{a}_\Omega) = -\bm{a}_\Omega$ 
(since $\bm{a}_\Omega \perp \FrameBasis$). 
Hence:
\begin{equation*}
\bm{e}^{(a)}(x) = \begin{pmatrix} (\FrameBasis \times \RotationAverage{\PlaneShape})\,\bm{a}_\Omega \\ 
-\FrameBasis \times \bm{a}_\Omega \end{pmatrix},
\end{equation*}
which is constant in $x$.

\paragraph{Torsion.}

The torsion mode displacement is
\begin{equation*}
\hat{\bm{u}}^{(\vartheta)} = \big( x\,(\FrameBasis \times \bm{r}) + \WarpTwistIesan\,\FrameBasis \big)\,a_\vartheta.
\end{equation*}
Applying the section-average operators:
\begin{align*}
\SectionAverage{\hat{\bm{u}}^{(\vartheta)}} &= x\,(\FrameBasis \times \CentroidLocation)\,a_\vartheta, \\
\RotationAverage{\hat{\bm{u}}^{(\vartheta)}} &= x\,\FrameBasis\,a_\vartheta + \RotationAverage{\WarpTwistIesan\,\FrameBasis}\,a_\vartheta 
= (x\,\FrameBasis + \CenterTwistLocation)\,a_\vartheta,
\end{align*}
where $\CenterTwistLocation := \RotationAverage{\WarpTwistIesan\,\FrameBasis} \in \mathbb{R}^3$ 
is the \emph{twist center} (note: $\CenterTwistLocation$ is generally not parallel to $\FrameBasis$).

Thus:
\begin{align*}
\bar{\bm{\theta}}^{(\vartheta)}(x) &= (x\,\FrameBasis + \CenterTwistLocation)\,a_\vartheta, \\
\bar{\bm{u}}^{(\vartheta)}(x) &= x\,(\FrameBasis \times \CentroidLocation)\,a_\vartheta 
- (x\,\FrameBasis + \CenterTwistLocation) \times \CentroidLocation\,a_\vartheta \\
&= x\,(\FrameBasis \times \CentroidLocation)\,a_\vartheta 
- x\,(\FrameBasis \times \CentroidLocation)\,a_\vartheta 
- (\CenterTwistLocation \times \CentroidLocation)\,a_\vartheta \\
&= -(\CenterTwistLocation \times \CentroidLocation)\,a_\vartheta.
\end{align*}
The strains are:
\begin{align*}
\bar{\bm{\kappa}}^{(\vartheta)} &= (\bar{\bm{\theta}}^{(\vartheta)})' = \FrameBasis\,a_\vartheta, \\
\bar{\bm{\gamma}}^{(\vartheta)} &= (\bar{\bm{u}}^{(\vartheta)})' + \FrameBasis \times \bar{\bm{\theta}}^{(\vartheta)} \\
&= \mathbf{0} + \FrameBasis \times (x\,\FrameBasis + \CenterTwistLocation)\,a_\vartheta \\
&= (\FrameBasis \times \CenterTwistLocation)\,a_\vartheta.
\end{align*}
Hence:
\begin{equation*}
\bm{e}^{(\vartheta)}(x) = \begin{pmatrix} \FrameBasis \times \CenterTwistLocation \\ \FrameBasis \end{pmatrix} a_\vartheta,
\end{equation*}
which is constant in $x$.

\paragraph{Flexure.}

The flexure mode displacement is
\begin{equation*}
\begin{aligned}
\hat{\bm{u}}^{(b)} &= (x\,\FrameBasis \times \bm{r} + \WarpTwistIesan\,\FrameBasis)\,c_\vartheta + \Big( -\tfrac{1}{6}x^3\,\mathbf{1}_\Omega 
+ x\,\PlaneShearShape 
+ \tfrac{1}{2}x^2\,\FrameBasis \otimes (\bm{r} - \CentroidLocation) 
+ \FrameBasis \otimes \WarpShearShape{} \Big)\,\bm{b}_{\Section},
\end{aligned}
\end{equation*}
where $c_\vartheta = \CenterShearLocation \cdot \bm{b}_{\Section}$ is the flexure-induced twist, 
and the shear center \(\CenterShearLocation\) is defined by \eqref{eq:def-shear-center}.


\subparagraph{Rotation.}
Applying $\RotationAverage{\cdot}$ to $\hat{\bm{u}}^{(b)}$:
\begin{align*}
\bar{\bm{\theta}}^{(b)}(x) &= (x\,\FrameBasis + \CenterTwistLocation)\,c_\vartheta 
- \tfrac{1}{2}x^2\,(\FrameBasis \times \bm{b}_{\Section}) 
+ x\,\RotationAverage{\PlaneShearShape}\,\bm{b}_{\Section} 
+ \RotationAverage{\FrameBasis \otimes \WarpShearShape{}}\,\bm{b}_{\Section} \\
&= (x\,\FrameBasis + \CenterTwistLocation)\,(\CenterShearLocation \cdot \bm{b}_{\Section}) 
- \tfrac{1}{2}x^2\,(\FrameBasis \times \bm{b}_{\Section}) 
+ x\,\RotationAverage{\PlaneShearShape}\,\bm{b}_{\Section} 
+ \RotationAverage{\FrameBasis \otimes \WarpShearShape{}}\,\bm{b}_{\Section}.
\end{align*}
where $\RotationAverage{\PlaneShearShape}\in\mathbb{R}^{3\times2}$, \(\RotationAverage{\FrameBasis\otimes\WarpShearShape{}}\in\mathbb{R}^{3\times2}\),
and \(\CenterTwistLocation := \RotationAverage{\FrameBasis\WarpTwistIesan}\in\Section\).
Then, grouping terms:
\begin{equation}
\bar{\bm{\theta}}^{(b)}(x) = -\tfrac{1}{2}x^2\,(\FrameBasis \times \bm{b}_{\Section}) 
+ x\,\left( \FrameBasis \otimes \CenterShearLocation + \RotationAverage{\PlaneShearShape} \right)\,\bm{b}_{\Section} 
+ \left( \CenterTwistLocation \otimes \CenterShearLocation + \RotationAverage{\FrameBasis\otimes\WarpShearShape{}} \right)\,\bm{b}_{\Section}.
\label{eq:cowper-theta-flexure}
\end{equation}

The curvature is:
\begin{equation}
\bar{\bm{\kappa}}^{(b)}(x) 
= -x\,(\FrameBasis \times \bm{b}_{\Section}) 
+ \left( \FrameBasis \otimes \CenterShearLocation + \RotationAverage{\PlaneShearShape} \right)\,\bm{b}_{\Section}.
\label{eq:kappa-flexure}
\end{equation}

\subparagraph{Deflection.}
Applying $\SectionAverage{\cdot}$ and subtracting $\bar{\bm{\theta}}^{(b)} \times \CentroidLocation$:
\begin{align*}
\SectionAverage{\hat{\bm{u}}^{(b)}} &= x\,(\FrameBasis \times \CentroidLocation)\,c_\vartheta 
- \tfrac{1}{6}x^3\,\bm{b}_{\Section} 
+ x\,\SectionAverage{\PlaneShearShape}\,\bm{b}_{\Section},
\end{align*}
using $\SectionAverage{\bm{r} - \CentroidLocation} = \mathbf{0}$ and $\SectionAverage{\WarpShearShape{}} = \mathbf{0}$.

Thus:
\begin{equation*}
\bar{\bm{u}}^{(b)}(x) = x\,(\FrameBasis \times \CentroidLocation)\,c_\vartheta 
- \tfrac{1}{6}x^3\,\bm{b}_{\Section} 
+ x\,\SectionAverage{\PlaneShearShape}\,\bm{b}_{\Section} 
- \bar{\bm{\theta}}^{(b)}(x) \times \CentroidLocation.
\end{equation*}

Computing $\bar{\bm{\theta}}^{(b)} \times \CentroidLocation$ from \eqref{eq:cowper-theta-flexure}:
\begin{align*}
\bar{\bm{\theta}}^{(b)} \times \CentroidLocation 
&= -\tfrac{1}{2}x^2\,(\FrameBasis \times \bm{b}_{\Section}) \times \CentroidLocation 
+ x\,\left( \FrameBasis \otimes \CenterShearLocation + \RotationAverage{\PlaneShearShape} \right)\,\bm{b}_{\Section} \times \CentroidLocation 
 + \big( \CenterTwistLocation \otimes \CenterShearLocation + \RotationAverage{\FrameBasis\otimes\WarpShearShape{}} \big)\,\bm{b}_{\Section} \times \CentroidLocation.
\end{align*}
Using $(\FrameBasis \times \bm{b}_{\Section}) \times \CentroidLocation = -\FrameBasis \otimes \CentroidLocation\,\bm{b}_{\Section}$:
\begin{equation*}
\bar{\bm{\theta}}^{(b)} \times \CentroidLocation 
= \tfrac{1}{2}x^2\,\FrameBasis \otimes \CentroidLocation\,\bm{b}_{\Section} 
+ x\,\bm{\Lambda}_1\,\bm{b}_{\Section} 
+ \bm{\Lambda}_0\,\bm{b}_{\Section},
\end{equation*}
where we define:
\begin{align}
\bm{\Lambda}_1 &:= (\FrameBasis \times \CentroidLocation) \otimes \CenterShearLocation 
+ \CentroidLocation^{\times} \RotationAverage{\PlaneShearShape}, 
\label{eq:def-Lambda1} \\
\bm{\Lambda}_0 &:= (\CentroidLocation \times \CenterTwistLocation) \otimes \CenterShearLocation 
+ \CentroidLocation^{\times} \RotationAverage{\FrameBasis\otimes\WarpShearShape{}}.
\label{eq:def-Lambda0}
\end{align}

Thus:
\begin{align*}
\bar{\bm{u}}^{(b)}(x) &= -\tfrac{1}{6}x^3\,\bm{b}_{\Section} 
- \tfrac{1}{2}x^2\,\FrameBasis \otimes \CentroidLocation\,\bm{b}_{\Section} 
+ x\,\left( (\FrameBasis \times \CentroidLocation) \otimes \CenterShearLocation 
+ \SectionAverage{\PlaneShearShape} - \bm{\Lambda}_1 \right)\,\bm{b}_{\Section} 
- \bm{\Lambda}_0\,\bm{b}_{\Section}.
\end{align*}

The derivative is:
\begin{equation*}
(\bar{\bm{u}}^{(b)})'(x) = -\tfrac{1}{2}x^2\,\bm{b}_{\Section} 
- x\,\FrameBasis \otimes \CentroidLocation\,\bm{b}_{\Section} 
+ \big( (\FrameBasis \times \CentroidLocation) \otimes \CenterShearLocation 
+ \SectionAverage{\PlaneShearShape} - \bm{\Lambda}_1 \big)\,\bm{b}_{\Section}.
\end{equation*}

The shear strain is:
\begin{equation*}
\bar{\bm{\gamma}}^{(b)}(x) = (\bar{\bm{u}}^{(b)})' + \FrameBasis \times \bar{\bm{\theta}}^{(b)}.
\end{equation*}

Computing $\FrameBasis \times \bar{\bm{\theta}}^{(b)}$ from \eqref{eq:cowper-theta-flexure}:
\begin{align*}
\FrameBasis \times \bar{\bm{\theta}}^{(b)} 
&= -\tfrac{1}{2}x^2\,\FrameBasis \times (\FrameBasis \times \bm{b}_{\Section}) 
+ x\,\FrameBasis \times \big( \FrameBasis \otimes \CenterShearLocation + \RotationAverage{\PlaneShearShape} \big)\,\bm{b}_{\Section} + \FrameBasis \times \big( \CenterTwistLocation \otimes \CenterShearLocation + \RotationAverage{\FrameBasis\otimes\WarpShearShape{}} \big)\,\bm{b}_{\Section} \\
&= \tfrac{1}{2}x^2\,\bm{b}_{\Section} 
+ x\,\FrameBasis \times \RotationAverage{\PlaneShearShape}\,\bm{b}_{\Section} 
+ \FrameBasis \times \big( \CenterTwistLocation \otimes \CenterShearLocation + \RotationAverage{\FrameBasis\otimes\WarpShearShape{}} \big)\,\bm{b}_{\Section},
\end{align*}
using $\FrameBasis \times (\FrameBasis \times \bm{b}_{\Section}) = -\bm{b}_{\Section}$ and 
$\FrameBasis \times (\FrameBasis \otimes \CenterShearLocation) = \mathbf{0}$ 
(since $\FrameBasis \times \FrameBasis = \mathbf{0}$).

Combining:
\begin{align*}
\bar{\bm{\gamma}}^{(b)}(x) &= \Big( -\tfrac{1}{2}x^2\,\bm{b}_{\Section} + \tfrac{1}{2}x^2\,\bm{b}_{\Section} \Big) \\
&\quad + x\,\Big( -\FrameBasis \otimes \CentroidLocation + \FrameBasis \times \RotationAverage{\PlaneShearShape} \Big)\,\bm{b}_{\Section} \\
&\quad + \Big( (\FrameBasis \times \CentroidLocation) \otimes \CenterShearLocation 
+ \SectionAverage{\PlaneShearShape} - \bm{\Lambda}_1 
+ \FrameBasis \times (\CenterTwistLocation \otimes \CenterShearLocation + \RotationAverage{\FrameBasis\otimes\WarpShearShape{}}) \Big)\,\bm{b}_{\Section}.
\end{align*}

The $x^2$ terms cancel. Define:
\begin{align}
\bm{\Gamma}_1 &:= -\FrameBasis \otimes \CentroidLocation , 
\label{eq:def-Gamma1} \\
\bm{\Gamma}_0 &:= (\FrameBasis \times \CentroidLocation) \otimes \CenterShearLocation 
+ \SectionAverage{\PlaneShearShape} - \bm{\Lambda}_1 
+ \FrameBasis \times (\CenterTwistLocation \otimes \CenterShearLocation + \RotationAverage{\FrameBasis\otimes\WarpShearShape{}}).
\label{eq:def-Gamma0}
\end{align}

Thus, because \(\FrameBasis \times \RotationAverage{\PlaneShearShape} = \mathbf{0}\):
\begin{equation}
\bar{\bm{\gamma}}^{(b)}(x) = x\,\bm{\Gamma}_1\,\bm{b}_{\Section} + \bm{\Gamma}_0\,\bm{b}_{\Section}.
\label{eq:gamma-flexure}
\end{equation}

\paragraph{Trace Matrices.}

Collecting the contributions from all four Ieşan modes, the beam strains are:
\begin{equation}
\bm{e}(x) = \TraceDeDpRoot\,\bm{p} + x\,\TraceDeDpOne\,\bm{p}
\end{equation}
where, partitioning $\bm{e} = (\epsilon, \bm{\gamma}_\perp, \varkappa, \bm{\kappa}_\perp)^\top$ 
and $\bm{p}$ is the vector of load parameters \eqref{eq:def-iesan-vector}:

\begin{equation}
\TraceDeDpRoot = \begin{pmatrix}
1 & \mathbf{0}^\top & 0 & (\bm{\Gamma}_0)_\parallel^\top \\[4pt]
\mathbf{0} & \mathbf{0} & (\FrameBasis \times \CenterTwistLocation)_\perp & (\bm{\Gamma}_0)_\perp \\[4pt]
0 & \mathbf{0}^\top & 1 & (\CenterShearLocation + (\RotationAverage{\PlaneShearShape})_\parallel)^\top \\[4pt]
\mathbf{0} & -\FrameBasis^\times & \mathbf{0} & (\RotationAverage{\FrameBasis\otimes\WarpShearShape{}} + \CenterTwistLocation \otimes \CenterShearLocation)_\perp
\end{pmatrix}
\label{eq:T-section-avg}
\end{equation}

For generic section geometry (non-degenerate warping functions), the six columns 
are linearly independent, giving $\mathrm{rank}(\mathbf{T}) = 6$ and hence the warp index is $w = 0$.

\begin{remark}[Simplification in centroidal coordinates]
When the coordinate origin coincides with the centroid ($\CentroidLocation = \mathbf{0}$), 
several terms simplify:
\begin{itemize}
\item $\bm{\Lambda}_1 = \CentroidLocation^\times \RotationAverage{\PlaneShearShape} = \mathbf{0}$
\item $\bm{\Lambda}_0 = \CentroidLocation^\times \RotationAverage{\FrameBasis\otimes\WarpShearShape{}} = \mathbf{0}$
\item $\bm{\Gamma}_1 = \FrameBasis \times \RotationAverage{\PlaneShearShape}$
\item $\bm{\Gamma}_0 = \SectionAverage{\PlaneShearShape} + \FrameBasis \times (\CenterTwistLocation \otimes \CenterShearLocation + \RotationAverage{\FrameBasis\otimes\WarpShearShape{}})$
\end{itemize}
For doubly-symmetric sections with centroidal coordinates, additional 
simplifications occur: $\CenterTwistLocation = \mathbf{0}$, $\CenterShearLocation = \mathbf{0}$, 
and many cross-coupling terms vanish.
\end{remark}

\begin{remark}[Comparison with classical formulas]
The tensor $\bm{\Gamma}_0$ in \eqref{eq:def-Gamma0}, when restricted to the 
transverse components and for centroidal coordinates, reduces to
\begin{equation}
(\bm{\Gamma}_0)_\perp = \SectionAverage{\PlaneShearShape} 
+ \FrameBasis \times \RotationAverage{\FrameBasis \otimes \WarpShearShape{}},
\end{equation}
which matches the classical shear-flexure coupling tensor 
$\bar{\bm{\Xi}}^{-1}$ of \eqref{eq:def-bfrombeta-avg}.
\end{remark}
